% Generated draft: requires mathematical and source review.
\documentclass[11pt]{article}
\usepackage[margin=25mm]{geometry}
\usepackage{fontspec}
\setmainfont{Latin Modern Roman}
\usepackage{amsmath,amssymb,mathtools,unicode-math}
\setmathfont{Latin Modern Math}
\usepackage{enumitem,needspace,xurl,hyperref}
\setlength{\parindent}{0pt}
\setlength{\parskip}{5pt}
\setlength{\emergencystretch}{3em}
\newcommand{\N}{\mathbb{N}}
\newcommand{\Z}{\mathbb{Z}}
\newcommand{\Q}{\mathbb{Q}}
\newcommand{\R}{\mathbb{R}}
\newcommand{\C}{\mathbb{C}}
\newcommand{\F}{\mathbb{F}}
\newcommand{\A}{\mathbb{A}}
\newcommand{\PP}{\mathbb{P}}
\newcommand{\E}{\mathbb{E}}
\newcommand{\eps}{\varepsilon}
\newcommand{\dd}{\,\mathrm{d}}
\begin{document}
\section*{UnsolvedMath ID: 10900026}
\textbf{Source problem code:} catalogue entry
\subsection{Definitions and mathematical statement}
Call a connected \(3\)-manifold hyperbolic if it admits a complete metric of sectional curvature \(-1\), with no finite-volume assumption. For \(H\leq\pi_1(M)\), let \(M_H\to M\) be the corresponding connected cover. The source defines \(M\) to be locally hyperbolic when every \(M_H\) with \(H\) finitely generated admits such a complete hyperbolic metric; it reports these covers as tame, meaning interiors of compact \(3\)-manifolds.

Agol asks to characterize hyperbolic \(3\)-manifolds with infinitely generated fundamental group. The singled-out existence question is whether there is a locally hyperbolic \(M\) whose fundamental group has no infinitely divisible subgroup such as \((\Q,+)\), although \(M\) itself admits no complete hyperbolic metric. Under the divisible-subgroup reading of the source terminology, a nontrivial subgroup \(D\) is excluded when
\[
 \forall h\in D\ \forall n\geq1\ \exists x\in D:\quad x^n=h.
\]
Thus the exclusion concerns nontrivial divisible subgroups, not merely a subgroup isomorphic to \(\Q\). The phrase infinitely divisible could instead refer to elements having roots in the ambient group; that convention is not specified in the excerpt and should be checked.

The supplied research summary reports that Tommaso Cremaschi constructed a locally hyperbolic manifold \(M_\infty\) with no divisible subgroup in its fundamental group that is not homeomorphic to any complete hyperbolic manifold, answering this particular existence question. It explicitly leaves the broader characterization as a larger program. This is a reported literature result, not a new verification or a claim that the full classification request is solved. Orientation, boundary and manifold-category conventions are not stated here.
\subsection{Short English statement}
Characterize complete hyperbolic \(3\)-manifolds with infinitely generated groups. The specific locally-hyperbolic but nonhyperbolic example, with no divisible subgroup, is reported solved by Cremaschi; the broader classification is separate.
\subsection{Sources}
\begin{itemize}
\item Agol, attributed item 5.1 in Delp, Hoffoss and Manning, \emph{Problems In Groups, Geometry, and Three-Manifolds} (2015), supplied PDF page 4. \url{https://arxiv.org/pdf/1512.04620#page=4}; original-status link \url{https://arxiv.org/abs/1512.04620}. Recorded access 2026-07-29; extraction source-tex checked against rendered PDF text, as reported in the record.
\item Tommaso Cremaschi, \emph{A locally hyperbolic 3-manifold that is not hyperbolic} (2017), \url{https://arxiv.org/abs/1711.11568}, supplied as answering the particular existence question.
\item Tommaso Cremaschi, \emph{Hyperbolization of infinite-type 3-manifolds}, \url{https://arxiv.org/abs/1904.11359}, supplied as related continuation.
\item Local UnsolvedMath registry, \texttt{data/unsolvedmath/problems.json}, permanent ID \(10900026\), exact source code \texttt{catalogue entry}; supplied solved/SOLVED-IN-LITERATURE label applies to the specific existence question as described in its summary.
\end{itemize}
\end{document}
